\pagestyle{empty}

\hfill

\vfill

\pdfbookmark[0]{Colophon}{colophon}
\section*{Colophon}

% ====================================================================================
% Le colophon classicthesis d'origine a ete retire: il sollicitait l'envoi d'une
% carte postale a Andre Miede et ne disait rien de ce document-ci. L'attribution
% du style, elle, est conservee ci-dessous: elle est due.
%
% Ce colophon a desormais deux fonctions distinctes. Les separer nettement.
% ====================================================================================

% ============================ TODO texte (~15 lignes) ==============================
% (1) COMMENT CE MEMOIRE A ETE FABRIQUE. Factuel, sans emphase.
%
%   - compose en \LaTeX{} avec le style classicthesis d'Andre Miede et Ivo
%     Pletikosic, inspire de « The Elements of Typographic Style » de Robert
%     Bringhurst. Garder l'url du depot; retirer la sollicitation de carte
%     postale;
%   - fontes employees (Palatino / URW Palladio pour le texte, les petites
%     capitales espacees pour les titres) --- verifier dans
%     classicthesis-config.tex avant d'affirmer;
%   - figures explicatives dessinees en Ti\emph{k}Z; figures de simulation
%     exportees en PGFPlots depuis matplotlib;
%   - LA CHAINE DE REPRODUCTIBILITE, qui est le point qui merite d'etre
%     detaille ici plutot qu'ailleurs: chaque figure de simulation est produite
%     par une experience versionnee sous \texttt{qanat}, qui enregistre le
%     commit, la graine et les parametres; l'export produit conjointement la
%     figure et son fichier de provenance, traduit en macros \LaTeX{} appelees
%     par la legende. La ligne grise sous chaque figure n'est donc pas recopiee
%     a la main et ne peut pas se desynchroniser. Renvoyer au depot public;
%   - bibliographie geree par biblatex/bibtex8.
% ====================================================================================

% ============================ TODO texte (~12 lignes) ==============================
% (2) DECLARATION D'USAGE D'OUTILS GENERATIFS (Claude, Anthropic).
%
% NE PAS SOUS-DECLARER: une formulation qui minimise est pire qu'une absence de
% mention, parce qu'elle invite le lecteur a faire l'inference lui-meme. La
% liste ci-dessous correspond a l'usage reel; l'ajuster si elle depasse ce que
% vous souhaitez declarer, mais en connaissance de cause.
%
%   USAGE REEL, par ordre d'importance decroissante:
%     - CODE DE SIMULATION: redaction des experiences \texttt{qanat}, des
%       algorithmes correspondants dans \texttt{hdrlib} (couts et gradients
%       riemanniens, minimiseurs, moyennes de Frechet) et de leurs tests;
%       detection et correction d'erreurs, dont plusieurs erreurs de signe et
%       de conjugaison presentes dans du code deja publie; banc de
%       non-regression contre l'implementation MATLAB de reference;
%     - STRUCTURE ET REDACTION: mise en place des squelettes de chapitre
%       (enonces, equations, algorithmes, blocs a rediger), figures Ti\emph{k}Z
%       explicatives, carte des dependances, et redaction de certaines parties
%       de l'etat des lieux a partir de mes propres notes et des sources citees;
%     - HOMOGENEISATION DES NOTATIONS: conversion des equations de plusieurs
%       articles vers la convention unique adoptee ici;
%     - CHAINE DE COMPILATION: systeme de figures, provenance, externalisation;
%     - RELECTURE.
%
%   CE QUI N'EN RELEVE PAS, et qu'il faut dire aussi clairement: les resultats,
%   les demonstrations, le choix de ce qui figure dans ce memoire et de ce qui
%   n'y figure pas, les orientations scientifiques, et la responsabilite
%   entiere du contenu.
%
%   COHERENCE: cette declaration doit s'accorder avec
%   \cref{ch:declaration} (FrontBackmatter/Declaration.tex) et avec ce que
%   l'ecole doctorale exige. Verifier les deux avant de figer la formulation.
% ====================================================================================

\bigskip
